% =====================================================================
%  Chapter 26 review -- the 14-station revision contract.
%  Revision-only: tables, one-liners, pointers back via \secref/\thmref/
%  \defref/\figref/\coderef/\probref.  No re-derivations.
%  Summit: CF 1292D "Chaotic V." (2700, L6) -- an infinite divisor tree (node
%  x>1 joins x/(lowest prime divisor of x); root 1; depth(x)=Omega(x)).  n mind
%  fragments sit at factorial nodes k_i!; choose an assembly node P minimising
%  the sum of path lengths to every fragment (repeats at one node counted).  The
%  marked set is the distinct factorials; the virtual tree is the trie of their
%  prime factors in DECREASING order (two nodes share a branch at depth d iff
%  their d+1 largest prime factors agree), and the answer is the WEIGHTED
%  CENTROID / 1-median cost on it.  Akikaze's official descent (group fragments
%  by current largest prime, step down the strict-majority branch, track the
%  delta) is the fold.  All four official samples hand-replayed: 5, 6, 6, 11.
%  Exports prob:c30-summit, lst:c30-kit, lst:c30-summit.
%  Mixed set is EXACTLY 10, all live-verified this session (names/ratings/tags
%  re-pulled from the CF problemset API), un-owned, sorted ASCENDING by rating;
%  each hint routes via the signature (def:c30-signature) or REJECTS the virtual
%  tree and names the true tool.  The famous virtual-tree problems 613D (Ch.27),
%  809E (Ch.25) and 1923E/915F are owned elsewhere and NEVER re-linked here.
% =====================================================================

\chaptersection{Chapter 26 Summary}

One sentence: \emph{a virtual tree is the smallest \(\lca\)-closed view of a fixed tree
that keeps a marked set of \(k\) nodes -- you sort them by \(\tin\), insert the LCAs of
adjacent pairs, and link the result with one monotone stack, obtaining a \(2k-1\)-node
tree whose compressed edges carry real path lengths; run a bottom-up DP on \emph{it} and
you recompute, in \(\bigO{k\log k}\), exactly what an \(\bigO{n}\) DP on the whole tree
would, because no discarded node lies on any marked-to-marked path
(\(\secref{c30-construction}\), \(\thmref{c30-fold}\)).} Chapter~16 handed you \(\lca\) and
promised this compression; this chapter cashes it in. Each row below is a pointer, not a proof.

\begin{center}
\renewcommand{\arraystretch}{1.25}
\begin{tabular}{@{}p{0.22\linewidth} p{0.52\linewidth} p{0.20\linewidth}@{}}
\toprule
\textbf{Result} & \textbf{Statement} & \textbf{Where} \\
\midrule
Virtual tree & \(A(S)=S\cup\{\lca(x,y)\}\) is the minimal ancestor-closed superset of a
marked set, an induced tree whose maximal unbranched chains are compressed to weighted edges. &
\defref{c30-vt} \\
Size bound & Sorting \(S\) by \(\tin\), the \(k-1\) \emph{adjacent-pair} LCAs already close the
set, and \(|A(S)|\le 2k-1\) -- linear in the marked set, independent of \(n\). & \thmref{c30-size} \\
Construction & sort by \(\tin\), add adjacent LCAs, dedup, then one monotone stack hangs each
node on its nearest kept ancestor; edge weight \(=\depth(v)-\depth(u)\). \(\bigO{k\log k}\). &
\coderef{c30-build} \\
Weighted fold & A bottom-up DP on the weighted virtual tree returns the identical value to the
honest \(\bigO{n}\) DP on the real tree; nodes off every marked-to-marked path are safely dropped. &
\thmref{c30-fold}, \coderef{c30-fold} \\
Pair-sum & \(\sum_{\{x,y\}\subseteq S}\dist(x,y)=\sum_{e=(u,v)}\mathit{w}(e)\cdot\mathit{cnt}(v)(k-\mathit{cnt}(v))\)
over virtual edges. & \thmref{c30-pairsum}, \coderef{c30-pairsum} \\
Steiner on a tree & The minimal connected subgraph spanning \(S\) is the expanded \(A(S)\): forced,
unique, size \(\sum_e\mathit{w}(e)\). Graph Steiner is a \emph{different} tool (bitmask DP). &
\secref{c30-apps}, \secref{c30-boundary} \\
The pick & Fixed tree + answer confined to a marked set + \(\sum k\ll n\) \(\Rightarrow\) virtual tree;
else tree DP / reroot, bitmask Steiner, link-cut, or HLD/Euler. & \defref{c30-signature}, \defref{c30-pick} \\
\bottomrule
\end{tabular}
\end{center}

\medskip
{\small\textbf{Algorithms from this chapter.}}
\begin{center}
\renewcommand{\arraystretch}{1.2}
\begin{tabular}{@{}p{0.36\linewidth} l l >{\raggedright\arraybackslash}p{0.20\linewidth}@{}}
\toprule
\textbf{Algorithm} & \textbf{Prep} & \textbf{Per run} & \textbf{Code} \\
\midrule
Build the virtual tree of a marked set & \(\bigO{n\log n}\) LCA & \(\bigO{k\log k}\) & \coderef{c30-build} \\
CF 2127E: one virtual tree per colour & \(\bigO{n\log n}\) & \(\bigO{n\log n}\) total & \coderef{c30-ancient} \\
Weighted fold on the compressed tree & \(\bigO{k\log k}\) & \(\bigO{k}\) & \coderef{c30-fold} \\
CF 2252F: per-colour VT + count DP + greedy & \(\bigO{n\log n}\) & \(\bigO{n\log n}\) total & \coderef{c30-spectral} \\
Pairwise distance sum over a set & \(\bigO{k\log k}\) & \(\bigO{k}\) & \coderef{c30-pairsum} \\
Recognition gate / boundary pick & -- & -- & \defref{c30-signature}, \defref{c30-pick} \\
Virtual-tree kit (one listing) & -- & -- & \coderef{c30-kit} \\
CF 1292D (Summit) -- VT + weighted centroid & \(\bigO{K^2}\) sieve & \(\bigO{n+\text{descend}}\) & \coderef{c30-summit} \\
\bottomrule
\end{tabular}
\end{center}

\chaptersection{Master Decision Flowchart}

Post-study ladder: walk top to bottom, stop at the first branch that matches. It decides
\emph{whether the answer is confined to a small marked set on a fixed tree} (the virtual
tree's only reason to exist) and, if so, \emph{which fold} to run on the compressed tree --
a weighted DP, a pairwise sum, a Steiner size, or a weighted centroid. Everything that
fails the top gate is handed to a full-tree or different engine (tree DP \chref{c10},
reroot \chref{c20}, HLD \chref{c16}, Euler \chref{c24}, link-cut \chref{c23}, bitmask
Steiner).

\begin{lstlisting}[language={}, basicstyle=\ttfamily\small, numbers=none, frame=single]
  Is the answer a quantity over a SMALL MARKED SET S (k << n) on a FIXED tree?
  |
  |-- The tree is EDITED between queries (link / cut)?
  |     +-- YES: NOT a virtual tree -- no static tin order survives an edit.
  |         Dynamic forest + online path/subtree value -> link-cut (Ch.25, c23).
  |
  |-- The answer depends on EVERY node's own value (global sum, per-node count,
  |   diameter, a value hanging in unmarked subtrees)?
  |     +-- YES: A(S) would DROP nodes the answer needs. NOT a virtual tree.
  |         Whole-tree DP (Ch.13, c10) or rerooting (Ch.15, c20).
  |
  |-- The terminals live on a GRAPH and may meet at NEW (non-terminal) nodes?
  |     +-- YES: the union-of-paths picture breaks. NOT a virtual tree.
  |         Bitmask Steiner DP (Dreyfus-Wagner, O(3^k n)); see 26.5, mixed 152E.
  |
  |-- Range / path / subtree queries on a FIXED tree where NOTHING shrinks with k?
  |     +-- HLD (Ch.22, c16) or Euler tour + Fenwick (Ch.5, c24). A(S) = T, no gain.
  |
  |-- Fixed tree, marked set of size k, answer routed through marked-to-marked
  |   paths, and O(k log k) beats O(n) per run: THIS chapter (def:c30-signature).
  |     First build A(S): sort S by tin, add adjacent-pair lcas, dedup, stack-link
  |     with edge weight depth[v]-depth[u].  |A(S)| <= 2k-1 (thm:c30-size).
  |     lst:c30-build.
  |
  |     |-- DP over the induced subtree (per colour / per subset)?  (26.2)
  |     |     run the Ch.13 fold on A(S); each virtual edge carries its length L and
  |     |     a subtree terminal count cnt[v] (thm:c30-fold).  Solved CF 2252F.
  |     |
  |     |-- sum of distances over ALL pairs of S?  (26.3)
  |     |     sum w(e) * cnt[v] * (k - cnt[v]) over virtual edges (thm:c30-pairsum).
  |     |     lst:c30-pairsum.
  |     |
  |     |-- minimum connected subgraph spanning S (Steiner ON A TREE)?  (26.3)
  |     |     it is the expanded A(S): size = sum w(e); count connected sets with a
  |     |     two-state fold.  Free because the tree forces the connectors.
  |     |
  |     |-- 1-median / weighted-centroid assembly on the marked set?  (Summit)
  |         the marked nodes are the factorials; their LCAs form a trie of descending
  |         prime factors; step down the strict-majority branch.  CF 1292D: lst:c30-summit.
  |
  |-- NONE of the above -- you wanted a full-tree or a different engine.  See 26.5
        (def:c30-pick): offline node-set queries -> Mo (Ch.27, c21); all-pairs distance
        distribution -> centroid (Ch.24, c18).
\end{lstlisting}

\chaptersection{Common Mistakes Checklist}

\begin{itemize}[nosep, leftmargin=1.4em]
  \item \textbf{Closing the set with \(\binom{k}{2}\) pairs instead of the \(k-1\) adjacent ones.}
        Sorting by \(\tin\) and taking only adjacent-pair LCAs already gives all of
        \(A(S)\setminus S\) (\(\thmref{c30-size}\)); iterating every pair is \(\bigO{k^2\log n}\)
        and \tle for \(k\) up to \(n\). \(\coderef{c30-build}\).
  \item \textbf{Forgetting the dedup after inserting LCAs.} An adjacent LCA can equal a marked node
        or another LCA; leaving duplicates breaks the stack invariant and double-counts a node.
        Sort by \(\tin\) again and \code{unique} before linking. \(\coderef{c30-build}\).
  \item \textbf{Rooting the virtual tree at an arbitrary marked node.} Root it at the \emph{global}
        LCA (of the first and last \(\tin\)-sorted marked node); a wrong root mis-orients edge
        weights and every \(\mathit{cnt}\). \(\secref{c30-construction}\), \(\thmref{c30-pairsum}\).
  \item \textbf{Setting a compressed edge's weight to \(1\) instead of \(L=\depth(v)-\depth(u)\).}
        A virtual edge stands for a whole swallowed chain; a unit weight silently under-counts
        every distance that uses it -- the classic \textbf{WA}. \(\defref{c30-edgeweight}\).
  \item \textbf{Seeding \(\mathit{cnt}(v)\) with the real subtree size.} \(\mathit{cnt}\) counts
        \emph{marked} nodes only; a virtual subtree with no marked node below never exists (it was
        pruned), so \(\mathit{cnt}(v)\ge1\). Mixing the two corrupts the pair-sum. \(\thmref{c30-pairsum}\).
  \item \textbf{Charging a per-vertex cost once per incident virtual edge.} A kept node shared by two
        compressed edges gets double-billed; charge vertices once (fold through \(\mathit{cnt}\)),
        edges through \(\mathit{w}\). \(\secref{c30-apps}\).
  \item \textbf{Re-walking the whole tree per subset.} The entire point is that the fold runs on
        \(\bigO{k}\) nodes (\(\thmref{c30-fold}\)); if your DFS still visits \(n\) nodes per colour,
        you built a virtual tree for nothing. \(\coderef{c30-fold}\).
  \item \textbf{Using a virtual tree when every node matters, or on a graph / moving tree.} That is
        the boundary: tree DP / reroot, bitmask Steiner, or link-cut (\(\defref{c30-pick}\),
        \(\secref{c30-boundary}\)).
\end{itemize}

\chaptersection{WA Clinic}

Verdicts are symptoms; the causes are the short list this chapter drilled. Diagnose top-down --
every leaf names a listing.

\begin{figure}[H]
\centering
\begin{tikzpicture}[diagram,
    v/.style={boxnode, fill=primaryblue!12, draw=primaryblue, text width=22mm, align=center},
    b/.style={boxnode, fill=difforange!12, draw=difforange, text width=12mm},
    l/.style={boxnode, fill=successgreen!10, draw=successgreen, text width=48mm,
              font=\scriptsize, align=left}]
  \node[v] (verdict) {\textbf{Verdict on the judge}};
  \node[b, below left=8mm and 5mm of verdict] (wa) {\textbf{WA}};
  \node[b, below right=8mm and 5mm of verdict] (tle) {\textbf{TLE / RE}};
  \node[l, below right=1mm and 2mm of wa.south] (wa1)
    {a compressed edge carries weight \(1\) not \(L=\depth(v)-\depth(u)\), so every distance through
     it is short -- the pair-sum / centroid is wrong (\(\defref{c30-edgeweight}\), \(\coderef{c30-build}\))};
  \node[l, below left=1mm and 2mm of wa.south] (wa2)
    {wrong root (a marked node instead of the global LCA) or a duplicate LCA left in, mis-orienting
     edges and \(\mathit{cnt}\) (\(\thmref{c30-size}\), \(\coderef{c30-build}\))};
  \node[l, below right=1mm and 2mm of tle.south] (tle1)
    {\(\binom{k}{2}\) LCA calls to close the set, or a per-subset DFS still walking all \(n\) nodes
     (\(\thmref{c30-size}\), \(\thmref{c30-fold}\))};
  \node[l, below left=1mm and 2mm of tle.south] (tle2)
    {the marked set is actually \(\Theta(n)\) so \(A(S)=T\) and the compression buys nothing, or the
     problem is graph Steiner wearing a tree costume (\(\secref{c30-boundary}\), \(\defref{c30-pick}\))};
  \draw[flow] (verdict) -- (wa);
  \draw[flow] (verdict) -- (tle);
  \draw[flow] (wa) -- (wa1);
  \draw[flow] (wa) -- (wa2);
  \draw[flow] (tle) -- (tle1);
  \draw[flow] (tle) -- (tle2);
\end{tikzpicture}
\caption*{WA: a compressed edge got unit weight instead of its real length, or the tree is rooted at a
marked node / carries a duplicate LCA, flipping \(\mathit{cnt}\). TLE or RE: closing the set with all
pairs or a per-subset DFS that still visits \(n\) nodes, or the marked set is really \(O(n)\) and a
virtual tree was the wrong engine.}
\end{figure}

\chaptersection{Quick Reference Card}

\begin{quickref}[Chapter 26 at a glance]
\textbf{A virtual tree compresses a fixed tree to a marked set \(S\) plus the LCAs of
\(\tin\)-adjacent nodes -- at most \(2k-1\) vertices whose compressed edges carry real
lengths -- so an \(\bigO{n}\) DP on the marked structure drops to \(\bigO{k\log k}\).}

\medskip
\textbf{Build (26.1).} Sort \(S\) by \(\tin\); add \(\lca(x_i,x_{i+1})\); sort+dedup; sweep a
stack, popping until the top is an ancestor of the current node, then hang it there with
weight \(\depth(v)-\depth(u)\). \(|A(S)|\le 2k-1\) (\(\thmref{c30-size}\), \(\coderef{c30-build}\)).

\medskip
\textbf{DP (26.2).} Run Ch.~13's bottom-up fold on \(A(S)\); each virtual edge contributes a
per-real-edge quantity times its length, and \(\mathit{cnt}(v)\) counts marked terminals below
\(v\) (\(\thmref{c30-fold}\), \(\defref{c30-edgeweight}\), \(\coderef{c30-fold}\)).

\medskip
\textbf{Applications (26.3).} Pairwise distance sum \(=\sum\mathit{w}(e)\cdot\mathit{cnt}(v)(k-\mathit{cnt}(v))\)
(\(\thmref{c30-pairsum}\), \(\coderef{c30-pairsum}\)); Steiner on a \emph{tree} \(=\) expanded \(A(S)\),
size \(\sum\mathit{w}(e)\), connected counts a two-state fold.

\medskip
\textbf{Recognition (26.4).} Three gates: fixed tree? answer marked-set-induced? \(\sum k\) small
enough that \(O(k\log k)\) beats \(O(n)\)? All yes \(\Rightarrow\) virtual tree (\(\defref{c30-signature}\)).

\medskip
\textbf{The boundary (26.5).} Every node matters \(\to\) tree DP / reroot (\(\chref{c10},\chref{c20}\));
graph terminals meet at new nodes \(\to\) bitmask Steiner; moving topology \(\to\) link-cut
(\(\chref{c23}\)); non-shrinking ranges \(\to\) HLD / Euler (\(\chref{c16},\chref{c24}\))
(\(\defref{c30-pick}\)).
\end{quickref}

\begin{lstlisting}[caption={Chapter 26 kit: the virtual-tree build (sort by tin, close under adjacent lca, dedup, monotone stack with depth-difference weights), the weighted fold, the pairwise distance sum, and a weighted-centroid descent -- one compact skeleton each, in one place. ASCII only, identifiers spelled without underscores.}, label={lst:c30-kit}]
// ===== BUILD the virtual tree (26.1) =====
// tin[], depth[] known; lca(u,v), isAncestor(u,v) from Ch.16
int vtSeen[MAXN]; vector<pair<int,int>> vtKids[MAXN]; // (child,L)
vector<int> buildVT(vector<int> S){
    sort(S.begin(),S.end(),[](int a,int b){return tin[a]<tin[b];});
    vector<int> N=S;
    for(int i=0;i+1<(int)S.size();i++) N.push_back(lca(S[i],S[i+1])); //k1
    sort(N.begin(),N.end(),[](int a,int b){return tin[a]<tin[b];});
    N.erase(unique(N.begin(),N.end()),N.end());         // |N|<=2k-1
    vector<int> st; st.push_back(N[0]);
    for(int i=1;i<(int)N.size();i++){ int v=N[i];
        while(!isAncestor(st.back(),v)) st.pop_back();
        int u=st.back(); addVEdge(u,v,depth[v]-depth[u]); // weight L
        st.push_back(v); }
    return N; }       // N[0] is the global lca

// ===== WEIGHTED FOLD (26.2) =====
ll cnt[MAXN], acc[MAXN];
void vtDfs(int v){ cnt[v]=vtSeen[v]; acc[v]=0;
    for(auto[u,L]:vtKids[v]){ vtDfs(u); cnt[v]+=cnt[u];
        acc[v]+=acc[u]+perEdge(cnt[u])*L; } }

// ===== PAIRWISE DISTANCE SUM (26.3) =====
ll pairSum(int v,int k){ ll res=0;
    for(auto[u,L]:vtKids[v])
        res += (ll)L*cnt[u]*(k-cnt[u]) + pairSum(u,k);
    return res; }

// ===== WEIGHTED CENTROID / 1-MEDIAN descent (Summit shape) =====
// walk from root into the child whose marked weight is a strict majority
int centroid(int root){ int v=root;
    while(true){ int best=-1,bw=0;
        for(auto[u,L]:vtKids[v]) if(cnt[u]*2>cnt[root] && cnt[u]>bw){bw=cnt[u];best=u;}
        if(best<0) return v; v=best; } }                 // then cost = sum w*dist(., centroid)
\end{lstlisting}

\chaptersection{Write-From-Memory Drills}

Five signature skeletons, each lifted from a listing already in the chapter. Cover the marked line,
reproduce from memory, then check the answer one line below.

\begin{lstlisting}[caption={Drill 1 -- close a marked set under lca (from lst:c30-build)}]
sort(S by tin);  vector<int> N = S;
for (int i=0; i+1 < (int)S.size(); i++)
    N.push_back( lca(S[i], S[i+1]) );     // k-1 ADJACENT pairs, not C(k,2)
sort(N by tin);
N.erase( unique(N.begin(), N.end()), N.end() );   // now N = A(S), |N| <= 2k-1
\end{lstlisting}
\noindent{\scriptsize Key: after sorting by \(\tin\), every \(\lca\) of marked nodes already appears
between two \emph{neighbours} in that order (\(\thmref{c30-size}\)); the second sort+unique removes the
duplicates a repeated LCA introduces (\(\coderef{c30-build}\)).}

\begin{lstlisting}[caption={Drill 2 -- the monotone-stack link (from lst:c30-build)}]
vector<int> st; st.push_back(N[0]);                    // N[0] = global lca = root
for (int i=1; i < (int)N.size(); i++){ int v = N[i];
    while (!isAncestor(st.back(), v)) st.pop_back();   // climb to nearest kept ancestor
    int u = st.back();
    vtKids[u].push_back({ v, depth[v] - depth[u] });    // compressed edge weight = L
    st.push_back(v);
}
\end{lstlisting}
\noindent{\scriptsize Key: the stack is the current kept root-path; popping until the top is an
ancestor finds \(v\)'s nearest ancestor in \(A(S)\) in amortised \(\bigO{1}\); the weight is the
\emph{depth difference}, the number of real edges the chain swallowed (\(\defref{c30-edgeweight}\)).}

\begin{lstlisting}[caption={Drill 3 -- the weighted fold (from lst:c30-fold)}]
void vtDfs(int v){
    cnt[v] = marked[v] ? 1 : 0;                         // count MARKED terminals
    acc[v] = 0;
    for (auto [u,L] : vtKids[v]){
        vtDfs(u);
        cnt[v] += cnt[u];
        acc[v] += acc[u] + perEdge(cnt[u]) * L;         // per-edge quantity x its real length
    }
}
\end{lstlisting}
\noindent{\scriptsize Key: identical to Ch.~13's subtree fold, except each child edge is multiplied by
its length \(L\) and \(\mathit{cnt}\) counts only marked terminals -- so a \(\bigO{n}\) DP runs in
\(\bigO{k}\) with the same answer (\(\thmref{c30-fold}\), \(\coderef{c30-fold}\)).}

\begin{lstlisting}[caption={Drill 4 -- pairwise distance sum over a set (from lst:c30-pairsum)}]
ll pairSum(int v, int k){ ll res = 0;
    for (auto [u,L] : vtKids[v])
        res += (ll)L * cnt[u] * (k - cnt[u])            // each of L edges splits cnt[u] | k-cnt[u]
               + pairSum(u, k);
    return res;
}
\end{lstlisting}
\noindent{\scriptsize Key: an edge is on a marked-to-marked path iff exactly one endpoint is below it;
\(L\) swallowed edges share the same split, so the term is \(L\cdot\mathit{cnt}\cdot(k-\mathit{cnt})\)
(\(\thmref{c30-pairsum}\), \(\coderef{c30-pairsum}\)).}

\begin{lstlisting}[caption={Drill 5 -- weighted-centroid / 1-median descent (from lst:c30-summit)}]
int centroid(int root){ int v = root;
    while (true){
        int best = -1, bw = 0;
        for (auto [u,L] : vtKids[v])
            if (cnt[u] * 2 > cnt[root] && cnt[u] > bw){ bw = cnt[u]; best = u; }
        if (best < 0) return v;                 // no strict-majority child -> we are the median
        v = best;
    }
}
\end{lstlisting}
\noindent{\scriptsize Key: stepping into a child that holds \(>\) half the total weight lowers the sum
of weighted distances; when no child does, the node is the weighted centroid minimising
\(\sum \mathit{weight}\cdot\dist\) -- the Summit's assembly node (\(\probref{c30-summit}\)).}

\chaptersection{Complete Problem Checklist}

Every problem this chapter offers -- two fully solved in-section, worked traces on the running
example \(T,S=\{4,8,6,10\}\), and one Summit on its own divisor-tree cases. Tick \emph{only} after an
unaided solve; a hinted attempt counts as half.

\begin{center}
\footnotesize
\setlength{\tabcolsep}{4pt}
\renewcommand{\arraystretch}{1.25}
\begin{tabular}{@{}c p{0.16\linewidth} l p{0.20\linewidth} >{\raggedright\arraybackslash}p{0.26\linewidth} c@{}}
\toprule
\textbf{\#} & \textbf{Problem} & \textbf{Source} & \textbf{Key technique} & \textbf{What to look for} & \textbf{Done?} \\
\midrule
\multicolumn{6}{l}{\textbf{Section 26.1 -- construction}} \\
W26.1 & stack sweep of \(A(\{4,8,6,10\})\) & worked & \lv{4} sort-by-tin + adjacent LCA & closed set \(1,2,4,8,3,6,10\); six weighted edges; \(\dist(8,10)=6\) & \(\square\) \\
W26.1b & Ancient Tree & CF 2127E & \lv{4} one VT per colour & branch node of \(2+\) colour VTs is forced cutie; star test \(\to\) cost \(0\) & \(\square\) \\
\midrule
\multicolumn{6}{l}{\textbf{Section 26.2 -- DP on the compressed tree}} \\
W26.2 & weighted fold equals the \(\bigO{n}\) DP & worked & \lv{5} edge weight + \(\mathit{cnt}\) & a length-\(L\) edge carries \(L\) copies of its flow; nodes off-paths dropped & \(\square\) \\
W26.2b & Spectral Components & CF 2252F & \lv{6} per-colour VT + count + greedy & virtual edge worth \(\min(m_{\text{in}},m_c-m_{\text{in}})\); test \(1\to\) \texttt{1 3 -1 -1 -1} & \(\square\) \\
\midrule
\multicolumn{6}{l}{\textbf{Section 26.3 -- applications (worked only)}} \\
W26.3 & pair-sum + Steiner size on \(S\) & worked & \lv{5} \(\sum w\cdot\mathit{cnt}(k-\mathit{cnt})\) & pair-sum \(=26\) over six pairs; Steiner \(=8\) edges & \(\square\) \\
\midrule
\multicolumn{6}{l}{\textbf{Section 26.4 -- recognition (worked only)}} \\
W26.4 & route five specs through the gate & worked & \lv{4} three-question signature & fixed / marked-induced / small \(\sum k\); reject graph Steiner & \(\square\) \\
\midrule
\multicolumn{6}{l}{\textbf{Section 26.5 -- the method boundary (worked only)}} \\
W26.5 & six specs a VT must hand away & worked & \lv{5} which engine & every-node \(\to\) DP; graph \(\to\) bitmask; moving \(\to\) LCT; range \(\to\) HLD & \(\square\) \\
\midrule
\multicolumn{6}{l}{\textbf{Summit (chapter end)}} \\
W26 & Chaotic V. & CF 1292D & \lv{6} divisor-tree VT + centroid & factorials closed by \(\tin\); descend majority branch; \(\to\) 5, 6, 6, 11 & \(\square\) \\
\bottomrule
\end{tabular}
\end{center}

\chaptersection{Recommended Study Order}

Four focused days. Day~1 internalises the build and CF 2127E; day~2 the weighted fold and CF
2252F; day~3 the two application families; day~4 recognition, the boundary, the Summit and the
mixed set.

\begin{center}
\small
\renewcommand{\arraystretch}{1.25}
\begin{tabular}{@{}llp{0.55\linewidth}@{}}
\toprule
\textbf{Day} & \textbf{Focus} & \textbf{Tasks} \\
\midrule
1 & Read + do 26.1 & hand-sweep the stack on \(S=\{4,8,6,10\}\) to the six weighted edges (\(\coderef{c30-build}\)); prove \(|A(S)|\le2k-1\) (\(\thmref{c30-size}\)); then CF 2127E, replaying the star test to cost \(0\) (\(\probref{c30-ancient}\)) \\
2 & Read + do 26.2 & write the weighted fold cold (\(\coderef{c30-fold}\)); then CF 2252F, replaying test 1 to \texttt{1 3 -1 -1 -1} (\(\probref{c30-spectral}\)) \\
3 & Read + do 26.3 & derive the pair-sum \(=26\) and Steiner \(=8\) by hand against the six direct distances (\(\thmref{c30-pairsum}\), \(\coderef{c30-pairsum}\)) \\
4 & Read + do 26.4, 26.5, Summit & place five specs on the gate (\(\defref{c30-signature}\)) and six on the pick (\(\defref{c30-pick}\)); then CF 1292D, replaying the descent to \(5,6,6,11\) (\(\probref{c30-summit}\)); finish with the first 3 mixed entries \\
\bottomrule
\end{tabular}
\end{center}

\chaptersection{Difficulty Progression}

Tiers are competency gates -- each row names the \emph{skill} the tier unlocks.

\begin{center}
\small
\renewcommand{\arraystretch}{1.3}
\begin{tabular}{@{}lp{0.32\linewidth}p{0.40\linewidth}@{}}
\toprule
\textbf{Tier} & \textbf{Problems} & \textbf{What passing this tier proves} \\
\midrule
\lv{3} \ \(\bigstar\bigstar\bigstar\)
  & close a marked set; name \(|A(S)|\)
  & you sort by \(\tin\) and take only the \(k-1\) adjacent LCAs, dedup included \\
\lv{4} \ \(\bigstar\bigstar\bigstar\bigstar\)
  & the stack build; CF 2127E (mixed at CF~2200)
  & you hang every node on its nearest kept ancestor with the right edge weight \\
\lv{5} \ \(\bigstar\bigstar\bigstar\bigstar\bigstar\)
  & weighted fold, pair-sum / Steiner, recognition, mixed at CF~2300--2500
  & you make a \(\bigO{k}\) virtual DP match the \(\bigO{n}\) honest one \\
\lv{6} \ \(\bigstar\bigstar\bigstar\bigstar\bigstar\bigstar\)
  & CF 2252F, CF 1292D (Summit), mixed at CF~2400--2700
  & you \emph{recognise} the marked-set signature (or reject it) and fold a weighted centroid \\
\bottomrule
\end{tabular}
\end{center}

\chaptersection{Reduction Gallery}

Every entry is ``one marked-set shape \(\to\) one virtual-tree move (or an honest hand-off).'' Each
line collapses a problem family onto a listing already in the chapter.

\begin{itemize}[nosep, leftmargin=1.4em]
  \item ``per colour / per query, an aggregate over the vertices of that set and the paths joining them''
        \(\to\) one virtual tree per set, Ch.~13 fold on it (\(\thmref{c30-fold}\), \(\coderef{c30-build}\)).
  \item ``is this vertex a forced branch point of colour \(d\)?'' (CF 2127E) \(\to\) it is an internal node
        of colour \(d\)'s virtual tree; count memberships (\(\probref{c30-ancient}\)).
  \item ``sum of distances from a set to a chosen connected component'' (CF 2252F) \(\to\) per-colour VT,
        edge worth \(\min(m_{\text{in}},m_c-m_{\text{in}})\), greedy \(k_c-1\) heaviest (\(\probref{c30-spectral}\)).
  \item ``sum of \(\dist\) over all pairs of a \(k\)-set'' \(\to\) per-edge separation \(\sum w\cdot\mathit{cnt}(k-\mathit{cnt})\)
        (\(\thmref{c30-pairsum}\), \(\coderef{c30-pairsum}\)).
  \item ``minimum connected subgraph spanning \(S\) \emph{on a tree}'' \(\to\) the expanded virtual tree,
        size \(\sum\mathit{w}(e)\) -- forced, no search (\(\secref{c30-apps}\)).
  \item ``place a facility minimising the weighted sum of distances to the marked nodes'' (CF 1292D)
        \(\to\) weighted centroid / 1-median on \(A(S)\) (\(\probref{c30-summit}\), \(\coderef{c30-summit}\)).
  \item ``the value is summed over \emph{every} node'' \(\to\) \emph{not} a virtual tree: full tree DP
        or rerooting (\(\defref{c30-pick}\), \(\secref{c30-boundary}\)).
  \item ``connect \(k\) terminals on a \emph{graph}'' \(\to\) \emph{not} a virtual tree: they meet
        at new nodes, so bitmask Steiner DP (\(\secref{c30-boundary}\), mixed CF 152E).
  \item ``path/subtree ranges on a fixed tree, nothing shrinks with \(k\)'' \(\to\) HLD / Euler+Fenwick
        (\(\chref{c16}\), \(\chref{c24}\)); ``edges are edited online'' \(\to\) link-cut (\(\chref{c23}\)).
\end{itemize}

\chaptersection{Mixed Problem Set}

Ten entries, each live-verified against the current Codeforces problem set (names, ratings and tags
re-pulled this session) and absent from every other roster in the book (the famous virtual-tree problems
CF 613D in Ch.~27 and CF 809E in Ch.~25 are owned elsewhere and not re-linked here). None says ``use a
virtual tree''; several deliberately \emph{do not} want one, and the drill is to run them through the
signature (\(\defref{c30-signature}\)) and say so. Each hint names the reduction (or the rejection) and
the Chapter~26 listing it leans on, ordered ascending by rating.

\begin{mixedset}
  \practice{Compressed Tree}{CF 1901E}
    {https://codeforces.com/problemset/problem/1901/E}{CF 2200 L5}
    {Tags \emph{dfs, dp, greedy, trees}. The name and shape -- a big tree, a small set of special
     vertices, an answer that routes through the paths joining them -- is the signature firing on all
     three questions: fixed tree, marked-set-induced answer, \(\sum k\) small (\(\defref{c30-signature}\)).
     Compress to \(A(S)\) and run the Ch.~13 fold on it (\(\coderef{c30-build}\), \(\thmref{c30-fold}\)).
     This is the genuine virtual-tree case of the set.}
  \practice{Diameter Intersections}{CF 2253E}
    {https://codeforces.com/problemset/problem/2253/E}{CF 2200 L5}
    {Tags \emph{brute force, dfs, fft, graphs, trees}. The answer is about the tree's \emph{diameter}
     and paths crossing it -- a global structural quantity, not a per-set fold. Gate S2 fails: the object
     is defined on all nodes, so nothing marked-set-induced shrinks (\(\defref{c30-pick}\)). Use this to
     practise REJECTING the virtual tree; the real tool is a diameter analysis (the \emph{fft} tag is a
     red herring for the hard sub-case). Contrast 26.5's first spec.}
  \practice{Tree Colorings}{CF 2112E}
    {https://codeforces.com/problemset/problem/2112/E}{CF 2200 L5}
    {Tags \emph{combinatorics, dfs, dp, math, number theory}. A counting problem on colourings with a
     combinatorics/number-theory core; the DP is over the whole tree's structure, not over a small
     marked subset's induced subtree. Ask the gate: does dropping the unmarked nodes lose information?
     If the count depends on every node's colour choice, it is a plain tree-DP, not a virtual tree
     (\(\defref{c30-signature}\), \(\secref{c30-boundary}\)).}
  \practice{Nightcrawler}{CF 2254G}
    {https://codeforces.com/problemset/problem/2254/G}{CF 2200 L5}
    {Tags \emph{data structures, dfs, dp, greedy, sortings}. A greedy/sorting tree problem with a DFS
     order -- a strong hint that \(\tin\) matters, but check whether the answer is confined to a small
     marked set (VT) or is a global greedy over all nodes (tree DP). If each query names a tiny set whose
     induced subtree you re-fold, it collapses to \(\bigO{k\log k}\) via \(\coderef{c30-build}\); if the
     greedy runs over the whole tree, it does not shrink and stays a Ch.~13 job (\(\defref{c30-pick}\)).}
  \practice{Build a Tree}{CF 2247E}
    {https://codeforces.com/problemset/problem/2247/E}{CF 2300 L5}
    {Tags \emph{constructive algorithms, trees, two pointers}. A constructive task: assign values so
     constraints hold. There is no marked set whose aggregate you compute -- you \emph{build} an
     assignment. Recognition reflex: a construction is rarely a virtual tree; solve by the greedy / two
     pointer the tags advertise. A deliberate rejection (\(\secref{c30-boundary}\)'s warning bullet 1:
     named vertices, but no set-induced sum to compress).}
  \practice{Dogeforces}{CF 1494D}
    {https://codeforces.com/problemset/problem/1494/D}{CF 2300 L5}
    {Tags \emph{constructive, dsu, divide and conquer, dfs}. You reconstruct a hidden tree from a
     distance-like matrix with DSU and divide-and-conquer -- the whole node set is the object, not a
     marked subset, and the tree is being \emph{built} not given-and-compressed. Reject the virtual tree
     (\(\defref{c30-signature}\): S1 has no fixed tree yet); reach for DSU + the reconstruction.}
  \practice{Again Trees\ldots (Easy Version)}{CF 2194F1}
    {https://codeforces.com/problemset/problem/2194/F1}{CF 2300 L6}
    {Tags \emph{bitmasks, dfs, dp, trees}. The \emph{bitmasks} tag is the tell: this is a
     subset-DP, exactly the shape 26.5 routes to bitmask DP rather than a virtual tree. On a \emph{tree}
     a small marked set would compress free (\(\secref{c30-apps}\)); when the state must range over
     \(2^k\) terminal subsets the answer does \emph{not} live on a single induced subtree, so run the
     bitmask fold. Distinguish it from CF 2127E where the set \emph{is} fixed per colour.}
  \practice{Spanning Tree Queries}{CF 1633E}
    {https://codeforces.com/problemset/problem/1633/E}{CF 2400 L5}
    {Tags \emph{binary search, dsu, graphs, greedy, sortings}. A spanning-forest / bottleneck question
     answered by sorting edges, DSU, and binary-searching the threshold -- a \emph{graph} MST-shape, not
     a fixed-tree marked-set fold. Gate S1 fails (the structure is a graph and you choose edges), so the
     virtual tree is the wrong lens; hand it to DSU + binary search (\(\defref{c30-pick}\)).}
  \practice{Garden}{CF 152E}
    {https://codeforces.com/problemset/problem/152/E}{CF 2500 L5}
    {Tags \emph{bitmasks, dp, graphs, trees}. The canonical Steiner-on-a-graph problem: connect a small
     set of terminals in a weighted \emph{graph}, and the cheapest connector can be a non-terminal. That
     is precisely the ``graph, new branch nodes'' rejection of \(\secref{c30-boundary}\) -- solve with the
     Dreyfus--Wagner bitmask Steiner DP \(O(3^k n)\), \emph{not} a virtual tree. On a plain \emph{tree} the
     same question would be free via the expanded \(A(S)\) (\(\secref{c30-apps}\)); the graph is the veto.}
  \practice{Tree Queries}{CF 1254D}
    {https://codeforces.com/problemset/problem/1254/D}{CF 2700 L5}
    {Tags \emph{data structures, probabilities, trees}. A probabilistic tree question solved with a
     data-structure sweep over a \emph{fixed} full tree -- the marked set is not the driver of the answer,
     and expectation runs over global structure. Nothing here is per-subset compression; treat it as the
     sharpest rejection at the top of the band (compare \(\secref{c30-boundary}\): a big tree and an
     aggregate over the whole thing, so \(A(S)\) would drop mass the expectation needs).}
\end{mixedset}



\chaptersection{The Summit -- CF 1292D Chaotic V.}

\begin{notebox}[Why this is the chapter top rope]
Every in-section problem here compressed a marked set on a tree you were \emph{given} as
\(n\) explicit nodes. The Summit hands you a tree you can never write down: the nodes
are all positive integers, \(x>1\) joins \(x/\!\text{lpf}(x)\) (lowest prime divisor), and
the fragments sit at factorial nodes \(k_i!\) with \(k_i\le 5000\) -- so the tree is
astronomically deep (\(\Omega(5000!)\approx 30{,}000\)) and its node labels overflow any
machine integer. The compression is the whole solution: only the \(\le 5000\) \emph{distinct}
factorials matter, their \(\lca\) are the common \emph{descending-prime-factor prefixes}, so
the virtual tree is that trie, and the answer -- minimise the sum of path lengths to an
assembly node \(P\) -- is a \textbf{weighted centroid / 1-median} on it
(\(\thmref{c30-fold}\) for the fold, \(\secref{c30-apps}\) for the distance currency, and the
Ch.~24 centroid idea for the descent). It is the honest summit, not the highest badge: the
hardest problem this chapter fully derives and hand-replays to all four judge outputs
(\(5,6,6,11\)).
\end{notebox}

\solvedproblem{c30-summit}{Chaotic V.}{CF 1292D}
  {https://codeforces.com/problemset/problem/1292/D}{CF 2700 L6}
  {implicit divisor tree (node \(x\!>\!1\) attaches to \(x/\!\text{lpf}(x)\), root \(1\),
   depth \(\depth(x)=\Omega(x)\)); the marked set is the distinct factorials \(k_i!\) with
   multiplicities; their LCAs are common prime-factor prefixes in DECREASING order, so the
   virtual tree is that trie; the objective is the WEIGHTED CENTROID / 1-median, reached by
   Akikaze's descent -- repeatedly step into the branch holding a strict majority of fragments,
   tracking the cost delta \(\mathit{cur}-=\mathit{majority}\), \(\mathit{cur}+= (n-\mathit{majority})\)}

\begin{problemstatement}[CF 1292D Chaotic V. (statement, constraints and all four samples verified against the live CF page)]
The A.R.C.\ Library core is an undirected graph of \emph{infinitely many} nodes labelled by
all positive integers \(1,2,3,\dots\). The node \(x\) (\(x>1\)) is directly connected to the
node \(x/f(x)\), where \(f(x)\) is the \textbf{lowest prime divisor} of \(x\). (This makes a
tree rooted at \(1\): repeatedly dividing by the lowest prime factor peels primes in
increasing order and reaches \(1\), so \(\depth(x)=\Omega(x)\), the number of prime factors of
\(x\) with multiplicity.)

Vanessa's mind is split into \(n\) fragments; the \(i\)-th sits at the node \(k_i!\) (a
factorial). Choose a node \(P\) (the assembly node) so that the \emph{total length of the
paths from every fragment to \(P\)} is as small as possible; if several fragments share a
node, that node's distance to \(P\) is counted once per fragment. Output that minimum total.

\textbf{Input.} Line~1: \(n\) (\(1\le n\le 10^6\)). Line~2: \(n\) integers \(k_1,\dots,k_n\)
(\(0\le k_i\le 5000\)); fragment \(i\) is at \(k_i!\). Note \(0!=1!=1\).

\textbf{Output.} One integer -- the minimum sum of path lengths to \(P\).

\textbf{Samples (all four official).} \(n{=}3,\ k{=}[2,1,4]\to \mathbf{5}\).
\(n{=}4,\ k{=}[3,1,4,4]\to\mathbf{6}\). \(n{=}4,\ k{=}[3,1,4,1]\to\mathbf{6}\).
\(n{=}5,\ k{=}[3,1,4,1,5]\to\mathbf{11}\). (The note draws the first \(24\) nodes and states
the first answer uses \(P=1\), the second uses \(P=6=3!\).)
\end{problemstatement}

\begin{attemptgate}[Attempt it yourself first -- \lv{6}, budget a real sitting]
The tree is infinite and the fragment nodes are factorials up to \(5000!\). You cannot build
it, cannot even label its nodes. What small set of nodes could the answer possibly depend on?
\peek{an optimal \(P\) lies on at least one fragment-to-root path, so it lies in the union of
the paths between the fragment nodes -- exactly the virtual tree of the marked set
\(\{k_i!\}\). How many \emph{distinct} factorials can there be, given \(k_i\le 5000\)?}
\nudge{collapse the \(n\) fragments to counts \(\mathit{cnt}[k]\) over the \(\le 5001\) distinct
factorials. The objective \(\sum_v \mathit{weight}(v)\cdot\dist(v,P)\) is minimised at the
\textbf{weighted centroid} of the virtual tree: from any node, moving to a child that holds
\(>\tfrac{n}{2}\) of the total weight strictly lowers the sum, so walk into the majority child
until none is. Distances are \(\Omega\)-differences: \(\dist(x,y)=\Omega(x)+\Omega(y)-2\Omega(\lca)\).}
\fullhint{Two nodes branch by their prime factors in DECREASING order: they share the branch at
depth \(d\) iff their \(d{+}1\) largest prime factors agree. Hence the virtual tree of the
factorials is a \emph{trie} of each \(k_i!\)'s prime factors sorted largest-first, and its edges
have length \(1\) (one prime peeled). Precompute \(\mathit{expo}[i][p]\) = exponent of prime \(p\)
in \(i!\) cumulatively (\(\bigO{K^2}\), \(K=5000\)); then \(\Omega(i!)=\sum_p\mathit{expo}[i][p]\).
Start at \(P=1\) with \(\mathit{cur}=\sum_i \mathit{cnt}[i]\cdot\Omega(i!)\). Let each active
\(i\)'s \emph{current largest remaining prime} be its branch key; group fragment counts by that
key, step \(P\) down into the prime \(p\) holding a strict majority (\(\mathit{best}\cdot 2>n\)),
updating \(\mathit{cur}\mathrel{-}=\mathit{best}\), \(\mathit{cur}\mathrel{+}= (n-\mathit{best})\),
and \(\mathit{ans}=\min(\mathit{ans},\mathit{cur})\); then peel one \(p\) off every node that took
that branch and discard the nodes that did not. Stop when no branch is a majority -- \(\mathit{ans}\)
is the 1-median cost. Read \(\mathit{cur}\) as \code{long long}.}
\end{attemptgate}

\solstep{Understanding}
The infinite divisor tree is a tree the problem never lets you touch: labels blow past any
integer type and depth reaches \(30{,}000\). The rescue is that the answer depends on only the
marked nodes -- the fragments at \(k_i!\) -- and the paths joining them, i.e.\ their virtual
tree (\(\defref{c30-vt}\)). Two facts make that tree explicit and small. First, collapsing
identical factorials to a weight \(\mathit{cnt}[k]\) leaves \(\le 5001\) marked nodes.
Second, the ancestor relation is pure prime-factor bookkeeping: node \(x\)'s path to \(1\) is
read off the prime factors of \(x\) in \emph{decreasing} order (peel the largest first when you
descend from the root), so \(x\) is an ancestor of \(y\) iff \(x\)'s decreasing-prime list is a
prefix of \(y\)'s, and \(\lca(x,y)\) is their common prefix. Thus the virtual tree of the
factorials is the \term{trie} of their decreasing-prime-factor sequences, every edge of length
\(1\), with a weight \(\mathit{cnt}\) parked at the factorial leaves.

The objective is then a textbook 1-median: \(\sum_v \mathit{weight}(v)\cdot\dist(v,P)\) over tree
nodes is minimised at the \textbf{weighted centroid} -- the unique node after which no child
subtree holds a strict majority of the total weight -- because shifting \(P\) across an edge into
a side holding \(t\) of \(n\) units changes the sum by \((n-t)-t=n-2t\), negative exactly when
\(t>n/2\). So the whole problem is ``build the (implicit) virtual tree, then walk to its weighted
centroid, tracking the cost.'' That is \(\thmref{c30-fold}\)'s principle (a \(\bigO{k}\) structure
replaces the \(\bigO{n}\) tree) fused with Chapter~24's centroid move.

\solstep{Key Observations}
\begin{observation}
\textbf{Depth is a prime count.} \(\depth(x)=\Omega(x)\) (prime factors with multiplicity), so
\(\depth(k!)=\Omega(k!)=\sum_{j=2}^{k}\Omega(j)\), precomputed cumulatively in \(\bigO{K^2}\) with a
smallest-prime sieve (\(\defref{c30-edgeweight}\)'s currency, here unit-length edges).
\end{observation}
\begin{observation}
\textbf{Factorials are NOT a chain.} \(4!=24\) is not an ancestor of \(5!=120\): the root-to-\(120\)
path is \(1,5,15,30,60,120\) (primes \(5,3,2,2,2\)) and never visits \(24\). The induced structure
is the trie of \emph{decreasing} prime-factor sequences, so you must group by common prefix, not
assume \(k!\mid (k{+}1)!\) gives tree ancestry. (\(\thmref{c30-size}\) still caps the virtual tree
at \(2k-1\) marked-plus-LCA nodes.)
\end{observation}
\begin{observation}
\textbf{The optimum is the weighted centroid.} For the sum of weighted distances the move rule is
local: cross an edge into a component with \(>n/2\) weight and the total drops by \(2t-n>0\).
Since \(P\) must lie in the marked Steiner skeleton, the descent is confined to the virtual tree,
and the accumulated \(\mathit{ans}\) over the walk is the minimum (\(\probref{c30-summit}\)).
\end{observation}

\solstep{Think It With Me}
\begin{thinkbox}
\deadend{Materialise the divisor tree up to \(\max k_i!\) and DFS for the centroid}
{The nodes are labelled by integers as large as \(5000!\) -- they do not fit in \code{long long},
and the depth is \(\Omega(5000!)\approx 3\times10^4\), so the tree has far more nodes than you can
store or reach. You can never write the tree down; only the \(\le5001\) marked factorials and
their prefix-LCAs are addressable. \textbf{TLE / overflow}.}
\deadend{Assume the factorials form a chain \(1!\mid2!\mid\cdots\) and run a line 1-median}
{\(k!\) divides \((k{+}1)!\) as integers, but divisibility is NOT ancestry in this tree: \(4!=24\)
is not on the path to \(5!=120\). Treating the marked set as a line puts the median at the wrong
node and over-counts every branch that is really a shared prime prefix. This is the seductive WA:
the answer is a \emph{centroid on a trie}, not a median on a path. \textbf{WA}.}
\deadend{Recompute each factorial's full factorisation and its LCA per pair, on the fly, all \(n\)}
{With \(n\le10^6\) fragments and \(\Omega\) up to thousands, doing \(O(\Omega)\) work per fragment
per step is \(\tle\). You must (i) collapse to \(\mathit{cnt}[k]\) over \(\le5001\) distinct
\(k\), and (ii) descend the majority branch by grouping \emph{distinct} factorials, never the
\(10^6\) raw fragments. \textbf{TLE}.}
\aha{Collapse to the \(\le5001\) distinct factorials with multiplicities, note the virtual tree is
the trie of their \emph{decreasing} prime-factor sequences (unit edges), and walk the weighted
centroid: from \(P=1\) with \(\mathit{cur}=\sum_k\mathit{cnt}[k]\,\Omega(k!)\), repeatedly step into
the branch whose current-largest-prime group is a strict majority, updating
\(\mathit{cur}\mathrel{-}=\mathit{best}\) and \(\mathit{cur}\mathrel{+}=n-\mathit{best}\), record
\(\min\), and peel one prime off the nodes that took that branch. The running minimum is the 1-median
cost.}
\end{thinkbox}

\solstep{The Approach Ladder}
\begin{center}\renewcommand{\arraystretch}{1.25}
\begin{tabular}{@{}c p{0.23\linewidth} p{0.29\linewidth} l c@{}}
\toprule
\textbf{\#} & \textbf{Approach} & \textbf{Idea} & \textbf{Time} & \textbf{Verdict}\\
\midrule
1 & \lv{2}\ build the tree to \(\max k_i!\), DFS & write every node & -- & \textcolor{warningred}{\ding{55}\ \textbf{RE}} \\
2 & \lv{3}\ line 1-median over \(1!,2!,\dots\) & assume factorials are a chain & \(\bigO{n}\) & \textcolor{warningred}{\ding{55}\ \textbf{WA}} \\
3 & \lv{4}\ pairwise-LCA virtual tree, raw \(n\) & right model, no collapsing & \(\bigO{n\Omega}\) & \tle \\
4 & \lv{6}\ \textbf{trie of decreasing primes + centroid descent} & collapse to \(\le K\) counts, walk majority & \(\bigO{K^2 + n}\) & \ac \\
\bottomrule
\end{tabular}\end{center}
Rung~4 is the Summit: rungs~1--3 have the right \emph{goal} but the wrong \emph{object} (an
impossible tree, a false chain, or an un-collapsed set); realising the virtual tree is a
\emph{decreasing-prime-factor trie} and that the optimum is its weighted centroid is the whole win.

\solstep{Solution}
Read \(n\) and collapse the fragments into counts \(\mathit{cnt}[k]\) for \(0\le k\le 5000\).
Precompute \(\mathit{expo}[i][p]\) = the exponent of prime \(p\) in \(i!\), by carrying the row
of \((i{-}1)!\) forward and adding the factorisation of \(i\) (a \(\bigO{K^2}\) sieve-style DP).
Initialise \(P\) at the root: \(\mathit{cur}=\mathit{ans}=\sum_i \mathit{cnt}[i]\sum_p\mathit{expo}[i][p]\)
(that is \(\sum_i\mathit{cnt}[i]\,\Omega(i!)\)), and set each active \(i\)'s pointer
\(\mathit{bestPD}[i]\) to the largest prime still present in \(i!\). Now descend: while some
pointer exceeds \(1\), group all fragment counts by their pointer prime, take the prime
\(\mathit{bestPrime}\) with the largest group \(\mathit{bestGroup}\); if \(\mathit{bestGroup}\cdot2\le n\)
or \(\mathit{bestPrime}=1\), stop (no strict majority). Otherwise move \(P\) one edge down into
that branch: \(\mathit{cur}=\mathit{cur}-\mathit{bestGroup}+(n-\mathit{bestGroup})\),
\(\mathit{ans}=\min(\mathit{ans},\mathit{cur})\); then every node whose pointer equalled
\(\mathit{bestPrime}\) drops one copy of that prime and reveals its next-largest, and every node
that did NOT take the branch is parked (\(\mathit{bestPD}=1\)). Print \(\mathit{ans}\).

\solstep{Proof of Correctness}
\begin{proofbox}
\textbf{(Virtual tree is the right object.)} An assembly node minimising the sum of distances lies
in the union of paths between marked nodes (otherwise slide it toward that union and every term
drops), so confining the search to the virtual tree of \(S=\{k_i!\}\) loses no optimum
(\(\thmref{c30-fold}\)). \textbf{(Structure of that tree.)} In the divisor tree the root-to-\(x\)
path is the sequence of partial products of \(x\)'s prime factors in decreasing order, hence
\(x\) is an ancestor of \(y\) iff \(x\)'s decreasing list is a prefix of \(y\)'s; so the virtual
tree of any node set is the trie of their decreasing-prime-factor lists with unit edges, and
\(\depth(x)=\Omega(x)\) makes the root-cost \(\sum_i\mathit{cnt}[i]\Omega(k_i!)\) exact.
\textbf{(Descent finds the centroid.)} Crossing an edge into the side with \(t\) of \(n\) weight
changes the sum by \(n-2t\), which is negative iff \(t>n/2\); the loop moves exactly when a
strict-majority branch exists and takes the (necessarily unique) such branch, so each step
strictly lowers the sum and no local step is skipped -- the walk halts at the unique weighted
centroid, the global minimiser. \textbf{(Cost tracking.)} \(\mathit{cur}\) is updated by the same
\(-t+(n-t)\) rule at each real edge, and \(\mathit{ans}\) records the running minimum, which is
the value at the final centroid. \textbf{(Grouping by largest prime = descending the trie.)}
Partitioning active nodes by \(\mathit{bestPD}\) and peeling the majority prime advances every
relevant node down one trie level and freezes the others, maintaining the invariant that
\(\mathit{bestPD}[i]\) is the next trie edge below \(P\) on \(i!\)'s path. Hence \(\mathit{ans}\)
equals the true minimum sum of path lengths. \(\Box\)
\end{proofbox}

\solstep{Implementation}
\begin{lstlisting}[caption={CF 1292D Summit: collapse fragments to counts over the distinct factorials, precompute prime exponents of each factorial, then descend the decreasing-prime-factor trie into the strict-majority branch (weighted centroid) tracking the cost delta. Faithful to Akikaze's official solution. ASCII only, identifiers without underscores.}, label={lst:c30-summit}]
#include <bits/stdc++.h>
using namespace std;
typedef long long ll;
const int MAXK = 5000;
int n, cnt[MAXK+1];
int expo[MAXK+1][MAXK+1];            // expo[i][p] = exponent of prime p in i!
int main(){
    scanf("%d", &n);
    for (int i=0;i<n;i++){ int k; scanf("%d",&k); cnt[k]++; }
    // cumulative prime exponents of i!
    for (int i=2;i<=MAXK;i++){
        for (int p=0;p<=MAXK;p++) expo[i][p] += expo[i-1][p];
        int t=i;
        for (int x=2; (ll)x*x<=t; x++) while(t%x==0){ expo[i][x]++; t/=x; }
        if (t>1) expo[i][t]++;
    }
    // P at root: cost = sum over fragments of Omega(i!) = sum of exponents
    ll ans=0, cur=0; int bestPD[MAXK+1];
    for (int i=0;i<=MAXK;i++) bestPD[i]=1;
    for (int i=1;i<=MAXK;i++){
        if (cnt[i]==0){ bestPD[i]=1; continue; }
        for (int p=1;p<=MAXK;p++){
            ans += (ll)expo[i][p]*cnt[i];
            cur += (ll)expo[i][p]*cnt[i];
            if (expo[i][p]!=0) bestPD[i]=p;      // largest prime still in i!
        }
    }
    // descend the trie into the strict-majority branch (weighted centroid)
    while (true){
        int mx=1; for (int i=0;i<=MAXK;i++) if (bestPD[i]>mx) mx=bestPD[i];
        if (mx<=1) break;
        int freq[MAXK+1]; memset(freq,0,sizeof freq);
        for (int i=0;i<=MAXK;i++) freq[bestPD[i]] += cnt[i];
        int bestPrime=1, bestGroup=freq[1];
        for (int p=2;p<=MAXK;p++) if (freq[p]>bestGroup){ bestGroup=freq[p]; bestPrime=p; }
        if (bestGroup*2 <= n) break;             // no branch holds a strict majority
        if (bestPrime==1) break;                 // majority is the parked/root group
        cur -= bestGroup; cur += (n - bestGroup);
        if (cur < ans) ans = cur;
        for (int i=1;i<=MAXK;i++){               // peel one bestPrime off the branch
            if (bestPD[i]!=bestPrime) bestPD[i]=1;      // others leave the descent
            if (bestPD[i]==1) continue;
            expo[i][bestPD[i]]--;
            while (bestPD[i]>1 && expo[i][bestPD[i]]==0) bestPD[i]--;
        }
    }
    printf("%lld\n", ans);
    return 0;
}
// Only the <=5001 distinct factorials are ever touched; the infinite tree is the trie of their
// prime factors in DECREASING order. Each descent step is one edge of that virtual tree.
\end{lstlisting}

\solstep{Trace on the Official Samples}
\begin{tracebox}[Sample 2: collapse, then descend to the centroid \(P=6\)]
Weights: \(\mathit{cnt}[3]{=}1\) (at \(3!{=}6\)), \(\mathit{cnt}[1]{=}1\) (at \(1!{=}1\)),
\(\mathit{cnt}[4]{=}2\) (at \(4!{=}24\)); \(n{=}4\), half \(=2\).
Depths \(\Omega\): \(\Omega(1!){=}0,\ \Omega(3!){=}2,\ \Omega(4!){=}4\). Root cost
\(\mathit{cur}=1{\cdot}2 + 1{\cdot}0 + 2{\cdot}4 = 10\), \(\mathit{ans}{=}10\).
Largest primes: \(3!{=}2{\cdot}3\to3\); \(4!{=}2^3{\cdot}3\to3\); \(1!\to\) parked.
\textbf{Step 1:} group by pointer: prime \(3\) holds \(\mathit{cnt}[3]{+}\mathit{cnt}[4]=3>n/2\).
Move into it: \(\mathit{cur}=10-3+(4-3)=8\), \(\mathit{ans}=8\) (\(P=3\)). Peel one \(3\):
\(3!\to2\) (\(\Omega\) left \(1\)), \(4!\to2\). Park node \(1\) (\(\mathit{bestPD}[1]{=}1\)).
\textbf{Step 2:} group by pointer: prime \(2\) holds \(\mathit{cnt}[3]{+}\mathit{cnt}[4]=3>n/2\).
Move into it: \(\mathit{cur}=8-3+1=6\), \(\mathit{ans}=6\) (\(P=6\)). Peel one \(2\): \(3!\to1\)
(parked), \(4!\to2\) (now largest prime \(2\) with one fewer exponent, still descending).
\textbf{Step 3:} group by pointer: prime \(2\) holds only \(\mathit{cnt}[4]=2\); \(2\cdot2\not>n\)
\(\Rightarrow\) \textbf{stop}. Assembly node \(P=6\), answer \(\boxed{6}\) -- matching the note
(\(P=6\), distances \(0,2,2,2\)). \(\Box\)
\end{tracebox}

\begin{tracebox}[The other three, by the same root-cost / majority rule]
\textbf{Sample 1} \(([2,1,4])\): \(\mathit{cnt}\) at \(2!{\to}1,1!{\to}1,4!{\to}1\); root cost
\(\Omega(2!){+}\Omega(1!){+}\Omega(4!)=1+0+4=5\). Largest primes \(2,(-),3\) tie at group size
\(1\), never a strict majority of \(n{=}3\) (need \(>1.5\Rightarrow2\)) \(\Rightarrow\) stop at
\(P=1\), answer \(\boxed{5}\).\\
\textbf{Sample 3} \(([3,1,4,1])\): \(\mathit{cnt}[3]{=}1,\mathit{cnt}[1]{=}2,\mathit{cnt}[4]{=}1\);
root cost \(2+0+4=6\). Prime \(3\) group \(=\mathit{cnt}[3]{+}\mathit{cnt}[4]=2\), not \(>2\)
\(\Rightarrow\) stop at \(P=1\), answer \(\boxed{6}\).\\
\textbf{Sample 4} \(([3,1,4,1,5])\): \(\mathit{cnt}[1]{=}2,\mathit{cnt}[3]{=}1,\mathit{cnt}[4]{=}1,\)
\(\mathit{cnt}[5]{=}1\); root cost \(0{+}2{+}4{+}5=11\). Prime groups: \(3\to\{3,4\}=2\),
\(5\to\{5\}=1\); max \(2\), not \(>2.5\) \(\Rightarrow\) stop at \(P=1\), answer \(\boxed{11}\).
All four match the judge. \(\Box\)
\end{tracebox}

\solstep{Complexity}
\begin{complexitybox}
Reading and collapsing to counts is \(\bigO{n}\). The prime-exponent DP over \(i!\) for
\(i\le K=5000\) is \(\bigO{K^2}\) time and \(\bigO{K^2}\) integers of table (\(\approx2.5\cdot10^7\),
about \(100\) MB, inside the 512 MB limit). The centroid descent runs at most one step per peeled
prime along the majority path -- \(\bigO{K^2}\) in the worst case but tiny in practice because it
halts the moment no branch is a strict majority -- with \(\bigO{K}\) work per step. All sums fit in
\code{long long} (\(n\le10^6\), \(\Omega(5000!)\approx3\times10^4\), product \(\le3\times10^{10}\)).
Total \(\bigO{n+K^2}\). \textbf{Space} \(\bigO{K^2}\). \emph{No node of the real tree is ever
built} -- the compression to the \(\le K\) distinct factorials' decreasing-prime trie is the entire
trick (\(\thmref{c30-fold}\)).
\end{complexitybox}

\begin{hardeningbox}
\begin{itemize}[nosep,leftmargin=1.4em]
  \item \textbf{Do not treat factorials as a chain.} \(4!\nmid_{\text{tree}}5!\) (24 is not on the
        path to 120); the ancestry is by \emph{decreasing-prime prefix}. Line-median logic is the
        Summit's signature \textbf{WA} (\(\probref{c30-summit}\)'s second dead end).
  \item \textbf{Park the non-majority nodes, and stop on ties.} A node that did not take the
        majority prime leaves the descent (\(\mathit{bestPD}=1\)); and the walk must halt when
        \(\mathit{bestGroup}\cdot 2\le n\) -- a group of \emph{exactly} half is not a strict
        majority, so moving would not lower the sum.
  \item \textbf{Collapse before you compute.} \(\mathit{cnt}[k]\) over \(\le5001\) distinct \(k\),
        never the \(10^6\) raw fragments; per-fragment factorisation is \tle (third dead end).
  \item \textbf{Use \code{long long} for the cost, int for exponents.} The running sum
        \(\mathit{cur}\) is up to \(n\cdot\Omega(K!)\approx3\times10^{10}\); only the exponent
        table fits in \code{int}. Overflow here silently corrupts the centroid comparison.
\end{itemize}
\end{hardeningbox}

\begin{recallbox}[Chapter 26 in ten seconds]
A virtual tree closes a marked set \(S\) under \(\lca\) with \(k-1\) adjacent-pair calls after a
\(\tin\)-sort, dedups to \(|A(S)|\le2k-1\), and links it with one monotone stack giving each edge
weight \(\depth(v)-\depth(u)\) (\(\thmref{c30-size}\), \(\coderef{c30-build}\)). Run Ch.~13's fold
on that compressed tree and you get the honest \(\bigO{n}\) answer in \(\bigO{k}\)
(\(\thmref{c30-fold}\), \(\coderef{c30-fold}\)); pairwise sums are \(\sum w\cdot\mathit{cnt}(k-\mathit{cnt})\)
and tree-Steiner is the expanded \(A(S)\) (\(\thmref{c30-pairsum}\), \(\coderef{c30-pairsum}\)).
Fixed-tree + marked-set-induced + small \(\sum k\) is the gate (\(\defref{c30-signature}\)); else
hand off to tree DP, bitmask Steiner, HLD or link-cut (\(\defref{c30-pick}\)). Summit CF 1292D
(\(\probref{c30-summit}\)) compresses an infinite divisor tree to a decreasing-prime-factor trie and
reads its weighted centroid.
\selftest{
  \item Why is sorting by \(\tin\) and adding only adjacent-pair LCAs enough to close \(S\)? \hfill\lv{4}
  \item What weight does a compressed edge \((u,v)\) carry, and what does \(\mathit{cnt}(v)\) count? \hfill\lv{5}
  \item State the pairwise-distance-sum formula and the running-example value. \hfill\lv{5}
  \item Give the two rejected shapes a virtual tree must hand away, and to what. \hfill\lv{5}
  \item In the Summit, why are the factorials a \emph{trie}, not a \emph{chain}, and where is the optimum? \hfill\lv{6}
}
\end{recallbox}

\chaptersection{Tool Kit Inventory}

The machine assembled through Chapter~26. Earlier chapters' rows are carried forward; this
chapter's are appended. Pointers only.

\begin{center}
\small
\renewcommand{\arraystretch}{1.25}
\begin{tabular}{@{}p{0.24\linewidth} l c p{0.22\linewidth} >{\raggedright\arraybackslash\scriptsize}p{0.17\linewidth}@{}}
\toprule
\textbf{Tool} & \textbf{Since} & \textbf{Lv} & \textbf{Trigger phrase} & \textbf{Where} \\
\midrule
\multicolumn{5}{l}{\itshape Carried forward} \\
LCA + \(\dist\), binary lifting & 16 & \lv{4} & ``the inner loop the compression calls'' & \chref{c11}, \thmref{c11-distf} \\
Euler \(\tin/\tout\), ancestor test & 5 & \lv{3} & ``the sort key and the stack predicate'' & \chref{c24} \\
Bottom-up subtree fold & 13 & \lv{4} & ``run the same DP on fewer nodes'' & \chref{c10} \\
Weighted centroid / 1-median & 24 & \lv{5} & ``minimise the sum of distances to a set'' & \chref{c18} \\
\midrule
\multicolumn{5}{l}{\itshape Chapter 26 -- new this chapter} \\
Virtual tree \(A(S)\) & 26.1 & \lv{4} & ``compress a fixed tree to a marked set'' & \defref{c30-vt}, \coderef{c30-build} \\
Adjacent-pair closure, \(|A(S)|\le2k-1\) & 26.1 & \lv{4} & ``\(\bigO{k}\) skeleton, independent of \(n\)'' & \thmref{c30-size} \\
Weighted fold on \(A(S)\) & 26.2 & \lv{5} & ``a \(\bigO{k}\) DP that equals the \(\bigO{n}\) one'' & \thmref{c30-fold}, \coderef{c30-fold} \\
Pairwise distance sum & 26.3 & \lv{5} & ``sum \(\dist\) over all pairs of a set'' & \thmref{c30-pairsum}, \coderef{c30-pairsum} \\
Steiner-on-a-tree & 26.3 & \lv{5} & ``minimum connected subgraph spanning marked nodes'' & \secref{c30-apps} \\
Recognition gate & 26.4 & \lv{4} & ``is this even a virtual tree?'' & \defref{c30-signature} \\
Method-boundary pick & 26.5 & \lv{5} & ``what does the compression throw away?'' & \defref{c30-pick} \\
Virtual-tree kit (one listing) & 26 & \lv{5} & ``all the skeletons in one place'' & \coderef{c30-kit} \\
Summit: divisor-tree VT + centroid & Summit & \lv{6} & ``infinite tree, factorial marks, assembly node'' & \coderef{c30-summit} \\
\bottomrule
\end{tabular}
\end{center}

\chaptersection{Related Topics}

\begin{itemize}[nosep, leftmargin=1.4em]
  \item \chref{c11} (Ch.~16, LCA) --- backward, \textbf{promise discharged}: Ch.~16 ended by
        promising that a virtual tree compresses a \(k\)-node set via sort-by-tin plus adjacent-pair
        LCAs, with \(\thmref{c11-distf}\) and lifting as the inner loop. \(\secref{c30-construction}\)
        is exactly that build -- it \emph{calls} \(\lca\) \(O(k)\) times and never re-derives it.
  \item \chref{c24} (Ch.~5, Euler Tour / \(\tin/\tout\)) --- backward: the DFS order is the sort key
        and the ancestor predicate of the stack; the flatten is reused, not re-taught.
  \item \chref{c10} (Ch.~13, Tree DP) --- backward: the bottom-up subtree fold is run verbatim on the
        compressed tree (\(\thmref{c30-fold}\)); 26.2 adds only the edge-weight and \(\mathit{cnt}\) carry.
  \item \chref{c23} (Ch.~25, Link-Cut Trees) --- contrast and boundary: link-cut \emph{moves} a forest,
        a virtual tree \emph{shrinks} a fixed one; 26.5 hands every ``edges change'' spec to \chref{c23}
        (\(\defref{c30-pick}\)). Ch.~25's mixed set forward-cited CF 809E to this chapter.
  \item \chref{c18} (Ch.~24, Centroid Decomposition) --- forward/tool: the Summit's assembly node is a
        \emph{weighted centroid} (1-median), the same minimiser Ch.~24 uses for distance sums, now found
        on the virtual tree rather than the full one.
  \item \chref{c16} (Ch.~22, HLD) / \chref{c24} (Ch.~5) --- contrast: path/subtree \emph{ranges} on a
        fixed tree do not shrink with \(k\), so they stay with HLD/Euler; a virtual tree is for a
        marked set, not an interval (\(\secref{c30-boundary}\)).
  \item \chref{c21} (Ch.~27, Mo's Algorithm) --- forward: an offline batch of overlapping node-set
        queries with no inverse is Mo-on-trees, complementing the per-set virtual-tree fold
        (\(\secref{c30-boundary}\)). \textbf{This is the Next chapter.}
  \item Bitmask Steiner DP (outside Part~V) --- boundary: connecting terminals that may meet at new
        nodes is Dreyfus--Wagner \(O(3^k n)\), not a virtual tree; mixed CF 152E is the honest contrast.
\end{itemize}

\begin{lastword}[Last Word --- Chapter 26]
\emph{``A virtual tree is a fixed tree seen through a keyhole: sort a marked set by \(\tin\), add the
LCAs of adjacent pairs, and a monotone stack hands you a \(\bigO{k}\) skeleton whose compressed edges
remember their real lengths -- so every distance, count and connected structure among the marked nodes
survives the shrinking, and an \(\bigO{n}\) DP becomes \(\bigO{k\log k}\).''}

\medskip
Chapter~16 promised this compression and Chapter~26 paid it. Section 26.1 closed the marked set under
\(\lca\) and bounded it at \(2k-1\) (\(\thmref{c30-size}\)); 26.2 ran Chapter~13's fold on the shrunken
tree with the two carrying rules that make it faithful (\(\thmref{c30-fold}\)); 26.3 cashed it in on
the two workhorse families -- the pairwise sum \(\sum w\cdot\mathit{cnt}(k-\mathit{cnt})\) and the
forced tree-Steiner tree (\(\thmref{c30-pairsum}\)); 26.4 drilled the recognition gate so you reach for
the tool only when it fits (\(\defref{c30-signature}\)); and 26.5 drew the honest boundary where the
compression loses information and hands the work to tree DP, bitmask Steiner, HLD or link-cut
(\(\defref{c30-pick}\)). The Summit CF 1292D (\(\probref{c30-summit}\)) took the idea to its limit -- an
infinite tree you compress to a decreasing-prime-factor trie and solve with a weighted centroid.

\medskip
Carry three reflexes forward:
\begin{enumerate}[nosep, leftmargin=1.6em]
  \item \emph{Sort by \(\tin\), then close.} The \(k-1\) adjacent-pair LCAs are all you need
        (\(\thmref{c30-size}\)); never pay \(\binom{k}{2}\).
  \item \emph{Weight the edges by real length.} A compressed edge is worth \(\depth(v)-\depth(u)\);
        a unit weight is the quiet WA that breaks every distance downstream
        (\(\defref{c30-edgeweight}\)).
  \item \emph{Check the marked set is doing work.} If the answer needs every node, or the terminals
        sit on a graph and meet at new nodes, or the tree is being edited -- you do not have a virtual
        tree (\(\defref{c30-pick}\)).
\end{enumerate}

\medskip
And three habits: dedup after inserting LCAs so the stack invariant holds; root the virtual tree at the
global LCA, not a marked node; and price the run by \(\sum k\) versus \(n\) \emph{before} coding, so you
only build \(A(S)\) when it actually shrinks. You now own the one tree engine that answers by making a
big tree \emph{small} -- the static, compressive companion to every moving or path-based structure in Part~V.

\medskip
\textbf{Next:} \chref{c21} (Mo's Algorithm on Trees) turns the difficulty inside-out again. Where a
virtual tree re-derives the answer per marked set in \(\bigO{k\log k}\), Mo orders a large batch of
\emph{offline} queries so that consecutive sets differ by one node, amortising the cost to a single
add/remove step each -- the technique for when the queries overlap so heavily that rebuilding \(A(S)\)
every time wastes work Mo would share.
\end{lastword}
